\subsection{FAMILIES OF POSITIVE FINITE NUMBERS SUMMABLE IN R}

THEOREM 1. A family \((x_{t})\) of finite real numbers \(\geqslant 0\) is summable in \(\mathbf{R}\) if and only if the set of partial finite sums of the family is bounded above in \(\mathbf{R}\) . If so, the least upper bound of this set is the sum of the family \((x_{t})\) .

For each finite subset H of the index set I, let \(s_{H} = \sum_{i \in H} x_{i}\) ; since the \(x_{i}\) are \(\geqslant 0\) , the relation \(H \subset H'\) implies \(s_{H} \leqslant s_{H'}\) . In other words, the mapping \(H \to s_{H}\) is an increasing function on the directed set \(\mathfrak{F}(I)\) of finite subsets of I; therefore (§ 5, no. 2, Corollary to Theorem 2) it has a finite limit if and only if it is bounded above.

Remark. Let \((\mathrm{H}_{\lambda})\) be a family of finite subsets of I such that, for each finite subset H of I, there is an index \(\lambda\) such that \(H \subset H_{\lambda}\) ; then \((x_{t})\) is summable if and only if the family of the \(s_{H_{1}}\) is bounded above in R. In particular, let \((x_{n})\) be a sequence of finite real numbers \(\geqslant 0\) , and for each integer n let \(s_{n} = \sum_{p=0}^{n} x_{p}\) ; then the sequence \((x_{n})\) is summable in R if and only if, for one sequence of strictly increasing integers \((n_{k})\) , the partial sequence \((s_{n_{k}})\) is bounded above in R.

Examples. 1) For each number \(q\) such that \(0 \leqslant q < 1\) , the sequence \((q^n)\) ('geometric progression of ratio \(q''\) ) is summable in \(\mathbf{R}\) , since \(s_n = \frac{1 - q^{n+1}}{1 - q} \leqslant \frac{1}{1 - q}\) ; the sum of this sequence is \(\lim_{n \to \infty} s_n = \frac{1}{1 - q}\) .

\begin{enumerate}
\def\labelenumi{\arabic{enumi})}
\setcounter{enumi}{1}
\tightlist
\item
  Let \(a\) and \(b\) be two numbers such that \(0 \leqslant a < 1\) and \(0 \leqslant b < 1\) ; then the family \((a^{m}b^{n})_{(m,n) \in \mathbb{N} \times \mathbb{N}}\) is summable in \(\mathbb{R}\) . For each finite subset of \(\mathbf{N} \times \mathbf{N}\) is contained in a subset of the form \([0, p] \times [0, p]\) , and we have
\end{enumerate}

\[
\sum_ {m = 0} ^ {p} \sum_ {n = 0} ^ {p} a ^ {m} b ^ {n} = \left(\sum_ {m = 0} ^ {p} a ^ {m}\right) \left(\sum_ {n = 0} ^ {p} b ^ {n}\right) = \frac {1 - a ^ {p + 1}}{1 - a}. \frac {1 - b ^ {p + 1}}{1 - b} \leqslant \frac {1}{(1 - a) (1 - b)}.
\]

\begin{enumerate}
\def\labelenumi{\arabic{enumi})}
\setcounter{enumi}{2}
\tightlist
\item
  For each integer \(p > 1\) , the sequence \((n^{-p})\) ( \(n > 0\) ) is summable, since
\end{enumerate}

\[
s _ {2 ^ {n + 1}} - s _ {2 ^ {n}} = \sum_ {k = 1} ^ {2 ^ {n}} (2 ^ {n} + k) ^ {- p} <   2 ^ {n}. (2 ^ {n}) ^ {- p}
\]

and therefore, adding these inequalities together,

\[
s _ {2 ^ {n}} <   \frac {1}{1 - 2 ^ {1 - p}}.
\]

\begin{enumerate}
\def\labelenumi{\arabic{enumi})}
\setcounter{enumi}{3}
\tightlist
\item
The sequence \((1 / n)\) \((n > 0)\) is not summable in \(\mathbf{R}\) , since
\end{enumerate}

\[
s _ {2 ^ {n + 1}} - s _ {2 ^ {n}} = \sum_ {k = 1} ^ {2 ^ {n}} \frac {1}{2 ^ {n} + k} > \frac {2 ^ {n}}{2 ^ {n + 1}} = \frac {1}{2}
\]

and therefore, adding these inequalities together,

\[
s _ {2 ^ {n}} > n / 2
\]

so that the criterion of Theorem 1 is not satisfied.

\begin{enumerate}
\def\labelenumi{\arabic{enumi})}
\setcounter{enumi}{4}
\tightlist
\item
  Let \((\mathbf{I}_n)\) be a sequence of mutually disjoint non-empty open intervals, all contained in an interval of finite length \(l\) . The sum of the lengths of a finite number of intervals of this family is \(\leqslant l\) (§ 1, no. 5) and therefore the family of lengths of the \(\mathbf{I}_n\) is summable in \(\mathbf{R}\) , and its sum is \(\leqslant l\) .
\end{enumerate}

THEOREM 2 (Principle of comparison). Let \((x_{\iota})_{\iota\in I}\) and \((y_{\iota})_{\iota\in I}\) be two families of finite real numbers \(\geqslant 0\) , such that \(x_{\iota}\leqslant y_{\iota}\) for all \(\iota\) . If \((y_{\iota})\) is summable in \(\mathbb{R}\) then so is \((x_{\iota})\) , and we have \(\sum_{\iota}x_{\iota}\leqslant\sum_{\iota}y_{\iota}\) ; if in addition there is an index \(\chi\) such that \(x_{\chi}<y_{\chi}\) , then \(\sum_{\iota}x_{\iota}<\sum_{\iota}y_{\iota}\) .

The hypothesis implies that,

\[
\sum_ {\iota \in \mathbf {H}} x _ {\iota} \leqslant \sum_ {\iota \in \mathbf {H}} y _ {\iota},
\]

for every finite subset H of I and the first part of the theorem follows from this; the inequality relating the sums follows from the principle of extension of inequalities (§ 5, no. 2, Theorem 1). If \(x_{\chi} < y_{\chi}\) , then

\[
\sum_ {\iota} x _ {\iota} = x _ {\chi} + \sum_ {\iota \neq \chi} x _ {\iota} <   y _ {\chi} + \sum_ {i \neq x} y _ {\iota} = \sum_ {\iota} y _ {\iota}.
\]

This theorem provides the most commonly used criterion for deciding whether or not a sequence \((x_{n})\) of real numbers \(\geqslant0\) is summable in R: we try to compare the given sequence with a simpler sequence \((y_{n})\) of which we already know whether it is summable or not. If there exists a finite number a \textgreater{} 0 such that \(x_{n} \leqslant ay_{n}\) for all n from a certain point onwards, and if \((y_{n})\) is summable, then \((x_{n})\) is summable; if on the other hand there is a finite number b \textgreater{} 0 such that \(x_{n} \geqslant by_{n}\) for all n from a certain point onwards, and if \((y_{n})\) is not summable in R, then \((x_{n})\) is not summable in R. We shall see in a later volume how such comparison sequences may be obtained in the cases which arise most frequently.

Examples. 1) Let \(a\) be a finite real number \(>0\) , and consider the sequence \(\left(\frac{a^n}{n!}\right)\) ; let \(n_0\) be the smallest integer such that \(a < n_0\) . Then

for each \(n \geqslant n_0\) we have

\[
\frac {a ^ {n}}{n !} \leqslant \frac {a ^ {n _ {0}}}{n _ {0} !} \cdot \left(\frac {a}{n _ {0}}\right) ^ {n - n _ {0}};
\]

since \(q = \frac{a}{n_0} < 1\) , the sequence \((q^{n - n_0})\) is summable, and therefore so is the sequence \(\left(\frac{a^n}{n!}\right)\) .

\begin{enumerate}
\def\labelenumi{\arabic{enumi})}
\setcounter{enumi}{1}
\tightlist
\item
  Let \((a_{n})\) be a summable sequence of positive numbers. Since
\end{enumerate}

\[
\lim _ {n \rightarrow \infty} a _ {n} = 0,
\]

there exists an integer \(n_0\) such that \(a_n \leqslant 1\) whenever \(n \geqslant n_0\) . Consequently, for each \(n \geqslant n_0\) we have \(a_n^2 \leqslant a_n\) , and therefore the sequence \((a_n^2)\) is summable in \(\mathbb{R}\) . So is the sequence \((a_n^p)\) for each integer \(p > 1\) .

\begin{enumerate}
\def\labelenumi{\arabic{enumi})}
\setcounter{enumi}{2}
\tightlist
\item
  Let \(a\) and \(b\) be two numbers \(< 1\) ; then
\end{enumerate}

\[
\frac {1}{a ^ {m} + b ^ {n}} \leqslant \frac {1}{2 (\sqrt {a}) ^ {m} (\sqrt {b}) ^ {n}}
\]

and hence the family \(\left(\frac{1}{a^m + b^n}\right)\) is summable in \(\mathbf{R}\) .

COROLLARY. Let \((x_{i})_{i\in I}\) be a summable family of finite numbers \(\geqslant 0\) in R. If H is any subset of I, we have

\[
\sum_ {i \in \mathbf {H}} x _ {i} \leqslant \sum_ {i \in \mathbf {I}} x _ {i},
\]

and equality holds only if \(x_{i} = 0\) for all \(i \in \mathbb{C}\mathbb{H}\) .

